\documentclass[12pt]{article}
\usepackage[utf8]{inputenc}
\usepackage{amsmath, amssymb, amsthm}
\usepackage{mdframed}
\usepackage{geometry}

\geometry{a4paper, margin=1in}

\parskip=1ex
\title{Proofs with Sets}
\date{}

\begin{document}

\maketitle

If A, B, and C are sets, then $A \setminus (B \cup C) = (A \setminus B) \ () $

$(L \subset R)$
We have $L = A \setminus (B \cup C)$ adn $R = (A \setminus B) \cap (A \setminus C)$.

\begin{proof}
    By DeMorgan's Law, $x \notin B$ and $x \notin C$.

    Then, $x \in A \setminus B$ and $x \in A \setminus C$.

    Assume x$ \in R$. Then, $ x \in A \setminus B$ and $x \in A \setminus C$

    So $x \notin B \cup C$.

    It follows that $x \in A$ but $x \notin B \cup C$.
\end{proof}

Suppose A, B, and C are sets. Prove that if $B \subseteq C$, then $A \times B \subset A \times C$.

\begin{proof}
    Assume that $z \in A \times B$.

    Then, $x \in A$ and $y \in B$ such that $z = (x,y)$.

    Since $B \subseteq C$ and $y \in B$, it follows that $y \in C$.

    So $x \in A$ and $y \in C$.

    Then, $(x,y) \in A \times C$.

    So $z \in A \times C$.
\end{proof}

For any sets $A$ and $B$, $\mathbb{P}(A) \cap \mathbb{P}(B) = \mathbb{P}(A \cap B)$

\begin{proof}
    Assume $X \in P(A) \cap P(B)$.

    Then, $X \in P(A)$ and $X \in P(B)$

    So, $X \subseteq A$ and $X \subseteq B$.

    Then, $X \subseteq A \cap B$.

    Then $X \in P(A \cap B)$.

    For the converse, we assume $X \in P(A \cap B)$.

    Then, $X \subseteq A \cap B$. It follows that $X \subseteq A$ and $X \subseteq B$.

    So, $X \in P(A)$ and $X \in P(B)$.

    It follows that $X \in P(A) \cap P(B)$.
\end{proof}

$\forall A,B$, prove that $A \subseteq B \iff A \setminus B = \emptyset$.

\begin{proof}
    Assume $A \setminus B = \emptyset$.


    That is, $A \subset B$.


    For the converse, assume $A \setminus B \neq \emptyset$.
\end{proof}

\newpage{}

\begin{proof}
    Let $A = \{0, 1 ,2\}$, $B = \{2, 4\}$, and $C = \{ 3,5\}$.
    Then, $(B \cup C) = \{2, 3, 4, 5\}$. So $A \setminus (B \cup C) = {0, 1}$.
    Then, $(A \setminus B) = {0, 1}$ and $A \setminus C = \{ 0,1,2\}$
\end{proof}

\end{document}